\chapter{L'Effleureuse}
\sessioninfo{Session 4}{Mardi 8 septembre 2026}

Lyra vint les chercher avant le lever du soleil. Malgré le mois de juillet, l'air était étonnamment frais et une petite brise parcourait la vallée. Dans la pénombre qui précédait l'aube, Kael, Morrÿs, Obélik et Destos quittèrent le quartier général des Faiseurs de Veuves. Alan leur adressa un dernier signe tandis que le capitaine Virgile les regardait partir.

\illustration[0.9\textwidth]{arc-000-epilogue.png}
{\textit{Kael, Morrÿs, Obélik et Destos quittent Phandalin sous le regard d'Alan et du capitaine Virgile.}}
{depart_phandalin}

Lyra avait changé de tenue. Une cape sombre couvrait désormais ses épaules et un masque pâle, seulement percé de quelques ouvertures étroites, dissimulait entièrement son visage. Lorsqu'ils l'interrogèrent, elle leur expliqua que ce masque était lié à sa discipline monacale. Quant à savoir comment elle parvenait à voir à travers des trous aussi petits, sa réponse fut particulièrement simple : avec l'habitude, on s'y faisait.

\illustration[0.38\textwidth]{Lyra_Sunstricke_misericorde_pose.png}
{\textit{Lyra dans sa tenue monacale.}}
{lyra_misericorde}

La route jusqu'à Leilon devait leur prendre environ deux heures. Lyra ne cachait pas qu'elle aurait volontiers fait partie du voyage : Chult était, selon ses propres mots, \textit{the place to be}. C'était précisément pour cette raison que les Faiseurs de Veuves tenaient à rester discrets. Envoyer trop de membres importants de la guilde sur place aurait attiré les regards et, avec eux, des problèmes dont Yarno préférait se passer.

La marche leur donna l'occasion de mieux se connaître. Lyra avait vu les quatre aventuriers affronter ses techniques la veille et souhaitait comprendre d'où leur venait leur aptitude au combat. Ils n'étaient pas encore des combattants accomplis, mais ils s'étaient montrés suffisamment efficaces pour éveiller sa curiosité.

Kael parla de la caravane itinérante dans laquelle il avait vécu. Au fil de ses voyages, il avait côtoyé des soldats, des archers et quelques personnes capables d'employer la magie pour soigner. Chacun lui avait transmis une partie de ce qu'il savait. Kael s'était ainsi formé sur le terrain, rencontre après rencontre, avant de reprendre la route pour suivre sa propre voie.

Destos affirma pour sa part qu'il ne savait pas vraiment se battre. Il n'était, disait-il, qu'un vaisseau de la lumière. Il avait quitté son village avec un objectif simple : parcourir le monde afin d'y répandre cette lumière. Les sortilèges dont il disposait n'étaient à ses yeux qu'un moyen d'accomplir cette mission.

\begin{narrateur}
Un homme chargé d'apporter la lumière partout où il passait.

On aurait presque pu l'appeler Aziz.
\end{narrateur}

Obélik révéla avoir été fermier. Il n'avait reçu aucun entraînement particulier, mais travailler avec des bêtes exigeait de la force, de l'assurance et une certaine habitude des situations difficiles. Sa carrure faisait le reste.

Morrÿs expliqua enfin qu'il avait appris à lancer des couteaux la veille.

Le résultat était encore planté dans le plafond de la salle d'entraînement lorsque Obélik le lui avait récupéré.

Lyra ne sembla pas s'en inquiéter. Les talents du nain, lui assura-t-elle, seraient tout aussi importants que ceux des autres.

\scenebreak

Après deux heures de marche, le groupe atteignit Leilon. Un fameux trois-mâts, fin comme un oiseau, les attendait. Un homme leur fit signe depuis le pont. Il se nommait Corvis Quillfore et commandait le navire qui devait les conduire jusqu'à Chult.

\illustration[0.9\textwidth]{Corvis_Quillfore.png}
{\textit{Corvis Quillfore, capitaine de l'Effleureuse.}}
{corvis_quillfore}

Le bâtiment s'appelait \textit{l'Effleureuse}. Corvis en parlait comme d'un petit bijou et estimait qu'avec lui, la traversée durerait plus probablement vingt jours que quarante. Selon son capitaine, il surfait sur les vagues : un vrai bonheur. Morrÿs observa la coque et la qualité du travail réalisé sur le bois. Son instinct lui souffla rapidement que le navire n'était pas seulement bien conçu : quelque chose de magique l'habitait.

\illustration[0.9\textwidth]{l-effleureuse.png}
{\textit{L'Effleureuse, le fameux trois-mâts de Corvis Quillfore.}}
{effleureuse}

Corvis avait cependant un problème. Il n'avait été prévenu que la veille et ne disposait d'aucun équipage pour cette traversée. L'Effleureuse était faite pour fonctionner avec un équipage réduit, mais certainement pas avec son capitaine pour seul marin. Les quatre passagers allaient donc devoir apprendre en naviguant.

La répartition des postes tint compte des qualités de chacun. Obélik, le plus puissant du groupe, serait chargé de tirer sur les cordages. Morrÿs préférait rester près du pont plutôt que de grimper dans les hauteurs ; il assisterait donc Corvis à la navigation et prêterait main-forte au goliath. Kael, dont la vue était la plus perçante, rejoindrait la vigie en haut du mât.

Destos hérita de tout le reste.

Il devrait vérifier les lignes de vie, aider Obélik à manœuvrer les bouts et rejoindre le premier canon si la situation l'exigeait. Une fois à Chult, Corvis recruterait un véritable équipage avant de poursuivre sa route vers les îles Solaires. D'ici là, il devrait se contenter de quatre aventuriers qui n'avaient jamais navigué.

Les amarres furent larguées et les premières voiles hissées. L'Effleureuse quitta Leilon tandis que Kael prenait place dans la vigie. La traversée venait à peine de commencer lorsqu'une silhouette apparut à l'horizon. Kael alerta aussitôt le pont, sans pouvoir encore identifier le bâtiment qui approchait.

Corvis, lui, le reconnut.

C'était un navire de la garde de Waterdeep.

\illustration[0.9\textwidth]{blocus-waterdeep.png}
{\textit{Un navire de la garde de Waterdeep intercepte l'Effleureuse.}}
{blocus_waterdeep}

Le capitaine avait déjà subi plusieurs contrôles lors de ses précédents voyages et savait qu'il serait difficile d'y échapper. Bientôt, une voix porta au-dessus de l'eau : la Brigade d'Intervention Technique et Expéditive ordonnait à l'Effleureuse de s'arrêter pour permettre l'inspection de sa cargaison. C'était le fleuron de la garde de Waterdeep, et une équipe allait monter à bord.

\begin{narrateur}
Quand tout part en couille, appelez la BITE.
\end{narrateur}

Corvis prétendit rencontrer quelques difficultés de manœuvre afin de gagner du temps. Il conseilla aux fugitifs de trouver des vêtements et de se rendre aussi méconnaissables que possible. Kael se laissa glisser depuis le haut du mât et rejoignit ses compagnons dans la cabine. Obélik trouva rapidement de quoi se changer, puis aida les autres à faire de même.

\scenebreak

Le capitaine de la garde monta sur l'Effleureuse avec quatre soldats. Corvis l'accueillit en affirmant transporter des minerais vers les îles Solaires. L'explication parut d'abord suffisante, mais il restait encore à interroger les membres de cet équipage improvisé.

Le premier mensonge vint d'Obélik.

Il s'appelait Hervé.

Le capitaine le dévisagea. Un goliath de sa taille pouvait changer de vêtements, mais il lui était plus difficile de modifier son visage, et le nom choisi sur le moment ne rendait pas son histoire beaucoup plus convaincante. Destos s'approcha discrètement et employa un sortilège d'assistance. Avec cette aide providentielle, Hervé parvint à rendre son existence légèrement plus crédible.

Lorsque son tour arriva, Destos resta plus proche de la vérité. Il se présenta comme un prêtre itinérant venu du Nord afin de répandre la lumière, puis expliqua qu'il souhaitait gagner Chult. Le capitaine de la garde se retourna aussitôt vers Corvis : ne venait-il pas de parler des îles Solaires ? Destos se corrigea immédiatement, mais trop tard.

Le mot avait été prononcé.

Le garde s'intéressa ensuite à Kael et l'interpella en désignant sa peau bleue. Kael répondit que cette couleur venait probablement de son état de santé. Pour appuyer son explication, il se pencha brusquement au-dessus du bastingage et fit semblant de vomir dans la mer.

Le capitaine n'insista pas.

La maladie avait parfois ses avantages.

Morrÿs fut interrogé en dernier. Le capitaine tenait désormais une feuille sur laquelle figuraient quatre portraits : ceux de Kael, Obélik, Morrÿs et Alan. Il compara longuement le visage du nain avec celui dessiné sur l'avis de recherche.

La ressemblance était troublante.

Morrÿs ressemblait beaucoup à lui-même.

Obélik tenta alors de détourner l’attention en orientant la conversation vers le parcours du capitaine de la garde. Celui-ci expliqua qu’avant d’occuper ce poste, il avait travaillé comme marin dans l’import-export. Après avoir perdu sa licence, il avait accepté ce nouvel emploi parce qu’il fallait bien vivre. À l’entendre, contrôler des cargaisons n’était pas exactement la vocation dont il avait rêvé.

Cette confidence donna peut-être une idée au goliath. L’Effleureuse manquait justement d’équipage, et il commença presque à proposer au capitaine de les rejoindre. Corvis s’approcha lentement et l’interrompit avant que cette excellente idée ne produise des conséquences encore plus difficiles à gérer.

Destos tenta alors de détourner l'attention. Grâce à sa thaumaturgie, il provoqua un phénomène inquiétant sur le navire de la garde. Le capitaine appela Rico, puis ordonna à Muso de retourner voir ce qui se passait. Il regagna finalement son bâtiment avec trois de ses quatre soldats, ne laissant sur l'Effleureuse que le garde occupé à ligoter Morrÿs.

Destos continua d'entretenir la confusion autour des cordages. Obélik, lui, se rapprocha sans bruit. Morrÿs eut juste le temps de répandre une poignée de billes sur le pont avant d'être immobilisé. Le soldat glissa, s'écrasa au sol, puis se releva et termina de nouer ses liens.

Le capitaine de la garde, qui était retourné sur son bâtiment, revint chercher le prisonnier. Obélik plaça discrètement une main dans son dos. Il avait compris ce que Corvis préparait.

Le goliath trancha la corde qui reliait les deux navires.

\scenebreak

Tout bascula en un instant. Kael se jeta sur le capitaine de la garde et tenta de le projeter vers son propre bâtiment. Il parvint à le saisir, mais pas à le soulever suffisamment pour l'expédier par-dessus le bastingage.

Obélik s'attaqua au dernier soldat resté sur l'Effleureuse.

Lui ne manqua pas de force.

Le garde quitta le pont beaucoup plus vite qu'il n'y était monté et retomba sur son propre navire. Corvis tira brutalement sur le gouvernail. Les coques commencèrent à s'écarter pendant qu'il ordonnait à chacun de rejoindre son poste et de tirer sur les bouts.

Destos se précipita vers les cordages. Morrÿs, toujours ligoté, sortit son poignard et tenta de couper ses liens. La lame glissa, manqua la corde et s'enfonça dans son propre bras. Kael, lui, saisit une nouvelle fois le capitaine adverse.

Cette fois, il le projeta par-dessus bord.

L'homme disparut entre les deux navires dans une gerbe d'écume. Sur le bâtiment de la garde, un soldat prit son élan, courut vers le bastingage et sauta pour rejoindre l'Effleureuse avant que l'espace ne devienne infranchissable.

Son bond fut trop court.

Il rejoignit son capitaine dans l'eau.

La distance augmenta, mais la garde n'avait pas renoncé. Un arbalétrier ajusta Obélik et lui planta un carreau dans la jambe. Un second projectile fendit l'air, manqua le goliath et perfora un tonneau. Un liquide inconnu commença à s'en échapper tandis que l'Effleureuse prenait de la vitesse.

La poursuite venait de commencer.

Chacun rejoignit son poste. Obélik et Destos se chargèrent des cordages, Morrÿs assista Corvis et Kael remonta vers la vigie. Derrière eux, le navire de Waterdeep conservait leur sillage.

Destos fit gronder la magie sous la coque des poursuivants. De lourds craquements résonnèrent dans les profondeurs, comme si une créature gigantesque remontait vers la surface. La ruse sema suffisamment de confusion pour ralentir un instant la garde. Lorsqu'il recommença, personne ne fut dupé une deuxième fois.

Depuis les hauteurs, Kael surveillait la mer et la côte qui se rapprochait. Quelque chose apparut droit devant eux : des récifs et des hauts-fonds barraient leur route. Il alerta Corvis.

Le capitaine ne vira pas.

Il maintint son cap droit vers le danger.

\scenebreak

Obélik fut le seul à voir Corvis sortir un petit objet bleu de sa poche. Une faible lueur s'éveilla à sa surface. Aussitôt, le vent se renforça et gonfla les voiles de l'Effleureuse.

Le trois-mâts bondit.

La coque sembla s'arracher à l'eau. Le navire accéléra avec une violence soudaine et se mit à surfer sur les vagues, exactement comme Morrÿs l'avait pressenti en observant sa construction. Devant eux, les hauts-fonds grossissaient à chaque seconde.

Corvis cria à tout le monde de s'accrocher.

Cela allait secouer.

L'Effleureuse percuta les hauts-fonds. Le choc se propagea de la proue à la poupe dans un craquement assourdissant. Destos fut arraché au pont et précipité dans la cale, où sa tête heurta durement le plancher. Corvis bascula par-dessus le bastingage et disparut dans la mer.

Sa ligne de vie se tendit.

Le capitaine resta accroché à son navire, traîné dans son sillage. Obélik aperçut la corde, se précipita et ramena Corvis sur le pont. Il était inconscient.

Derrière eux, la garde tenta de suivre.

Son navire s'engagea à son tour au milieu des hauts-fonds. Il n'accéléra pas. Il ne se souleva pas. Sa coque s'enfonça dans le banc de sable et le bâtiment s'immobilisa brutalement, prisonnier de l'obstacle qu'il avait voulu franchir.

L'Effleureuse poursuivit sa course.

Ils étaient libres.

\scenebreak

Kael tenta de prodiguer des soins magiques à Corvis, sans parvenir à le réveiller. Morrÿs essaya à son tour, mais le capitaine demeura étourdi. Obélik décida alors d'employer une méthode plus directe.

Il le gifla.

\begin{narrateur}
On avait dit pas trop fort, Obélik.
\end{narrateur}

Corvis revint progressivement à lui. Son premier constat concerna la température de l'eau : elle était fraîche. Quelqu'un précisa aussitôt qu'à la nage, et elle était glacée.

Lorsqu'il fut en état de parler davantage, il expliqua que l'étrange pouvoir de l'Effleureuse ne pouvait pas être utilisé à volonté. Il faudrait attendre avant de recommencer. Il espérait surtout que la garde ignorerait leur destination réelle, même si Destos avait laissé échapper le nom de Chult pendant le contrôle.

Le clerc entreprit de vérifier qu'aucun dispositif de traçage n'avait été dissimulé à bord. Il lança un rituel de détection de la magie et parcourut le navire. L'expérience se révéla d'une efficacité toute relative : de la proue à la poupe, l'Effleureuse rayonnait de magie, mais aucun élément distinct ne permit de révéler la présence d'un traqueur.

Le danger écarté pour le moment, Kael remarqua une marmite. Il s'en approcha, inspira profondément et reconnut aussitôt l'odeur du ragoût de chevreuil servi la veille au quartier général des Faiseurs de Veuves. C'était exactement le plat qu'Alan leur avait préparé.

Corvis lui apprit que son frère avait cuisiné le repas.

Une question venait peut-être de trouver sa réponse. D'autres, en revanche, ne tardèrent pas à surgir. Morrÿs demanda s'il restait de la mousse au chocolat.

Corvis ne connaissait pas le chocolat.

Décidément, il y avait chez les Faiseurs de Veuves de sérieuses lacunes culinaires.
